\documentclass[a4paper,11pt]{article}
\usepackage[utf8]{inputenc}
\usepackage{amsmath,amssymb}
\usepackage{graphicx}
\usepackage{hyperref}
\usepackage{booktabs}
\usepackage[margin=1in]{geometry}
\usepackage{natbib}
\bibliographystyle{unsrtnat}

\newcommand{\fnl}{f_{\rm NL}}
\newcommand{\Mpl}{M_{\rm Pl}}

\title{Structural Barriers to Geometric Dark Energy and the\\Perturbation-Transparency of Einstein-Cartan-Holst Gravity}

\author{Houston Golden\\[4pt]
{\small\textit{Independent Researcher}}\\[2pt]
{\small\texttt{houston@hubify.com}}}

\date{\today}

\begin{document}
\sloppy
\maketitle

\begin{abstract}
We present a systematic catalog of structural barriers that prevent nonsingular bouncing cosmologies from generating late-time dark energy through first-principles mechanisms. Through 7 foundation studies and 17 research branches testing every viable theoretical route within the Einstein-Cartan-Holst (ECH) framework, we identify 14 independent barriers that close all standard mechanism classes. The barriers range from mass-coupling locks and topological-shift dualities to scale-separation arguments and perturbation-transparency theorems. Our central theoretical result is a perturbation-transparency theorem for minimal ECH gravity: for canonical scalar field matter, the spin density vanishes identically, torsion is algebraically zero at all perturbation orders, the Holst term reduces to the Pontryagin topological density, and the Barbero-Immirzi parameter becomes invisible in all scalar and tensor perturbation observables. This establishes ECH as a perturbation-transparent UV regulator---it resolves the Big Bang singularity through the spin-torsion modification of the Friedmann equation at Planck density, but leaves zero fingerprint on the primordial perturbations that survive into the CMB. We also document the systematic rejection of the hybrid dark-energy loophole: appending late-time dynamical-dark-energy freedom (e.g., CPL parametrization) to a bounce model was explored across 7 disguised forms and rejected on first-principles grounds, as it does not derive the dark-energy signal from the bounce mechanism. The observable science case for bouncing cosmology therefore rests entirely on generic, mechanism-independent predictions such as the matter-bounce non-Gaussianity $\fnl = -35/8$.
\end{abstract}

\section{Introduction}

The Einstein-Cartan-Holst (ECH) formulation of gravity extends general relativity by allowing spacetime torsion, sourced by the spin density of matter~\citep{Holst:1995pc}. In cosmological settings, the spin-torsion coupling produces a repulsive effective potential at Planck-scale densities, resolving the Big Bang singularity through a nonsingular bounce~\citep{Alexander:2014eva}. A natural question is whether this modified gravitational dynamics can also address other open problems in cosmology---particularly late-time cosmic acceleration (dark energy) and distinctive perturbation-level observables.

This paper presents the results of a systematic investigation of these questions, documenting both the negative structural results (14 barriers closing all standard dark-energy derivation routes) and the central positive theoretical finding (the perturbation-transparency of minimal ECH gravity).

\section{The 14 Structural Barriers}

We tested 7 foundation mechanism classes (Foundations A--G) and 7 additional observational channels (Branches H--O, plus the ECH perturbation gates) for the possibility of connecting the bounce to late-time dark energy or producing distinctive observable signatures. Each test yielded a named structural barrier.

\begin{table}[h]
\centering
\small
\begin{tabular}{clll}
\toprule
\# & Barrier & Source & Mechanism Blocked \\
\midrule
1 & Mass-Coupling Lock & Found.\ A & Propagating torsion as DE \\
2 & Topological-Shift Duality & Found.\ B & Geometric pseudoscalar mass protection \\
3 & Scalar-Tensor Universality & Found.\ C & Distinctive geometric content on FRW \\
4 & Planck Suppression & Found.\ D & Disformal / connection coupling effects \\
5 & Scale Separation & Found.\ E & Global vacuum integral coupling \\
6 & Attractor-Sensitivity Dilemma & Found.\ F & Initial-condition transfer to DE \\
7 & Parameter Immunity & Found.\ G & Cyclic vacuum selection \\
8 & Parity-Even Interaction & Branch H & Tensor chirality from the bounce \\
9 & Liouville Conservation & Branch J & Reversible state selection \\
10 & UV$\to$IR Specificity Dilemma & Branch L & Generic vs.\ bounce-specific bridge \\
11 & Decoupling Universality & Branch L/M & Light gauge field coupling \\
12 & Vacuum Amplification Ceiling & Branch M & Gravitational wave amplitude \\
13 & Gravitational Democracy & Branch N/O & Relics, baryogenesis, vacuum transitions \\
14 & Perturbation Transparency & ECH Gates & ECH-specific perturbation signatures \\
\bottomrule
\end{tabular}
\caption{The 14 structural barriers. Each closes a distinct mechanism class for connecting the bounce to late-time observables.}
\label{tab:barriers}
\end{table}

\subsection{Summary of Key Barrier Arguments}

\textbf{Barriers 1--4 (Foundations A--D):} These establish that single-field geometric dark energy from torsion is structurally impossible in the minimal framework. Foundation A shows that propagating torsion modes face a mass-coupling lock: the effective coupling $g_{\rm eff} \sim 1/(\Mpl\sqrt{|t_3|})$ requires graviton-loop fine-tuning at the level of 1 in $10^{57}$. Foundation B discovers a topological-shift duality: mass protection and geometric content cannot coexist simultaneously. Foundation C proves scalar-tensor universality on FRW backgrounds ($T_0 = Q_0 = 0$). Foundation D shows Planck suppression: connection couplings give 1 $\partial\phi$/vertex where 2 are needed for disformal effects.

\textbf{Barriers 5--7 (Foundations E--G):} These close non-minimal and global routes. Foundation E shows scale separation ($V_4^{\rm bounce}/V_4^{\rm total} \sim 10^{-60}$). Foundation F identifies the attractor-sensitivity dilemma. Foundation G establishes parameter immunity of the vacuum energy to bounce-scale physics.

\textbf{Barriers 8--13 (Branches H--O):} These close observational channels. Barrier 8 proves the spin-torsion effective interaction $(J^5)^2$ is parity-even. Barrier 12 establishes the vacuum amplification ceiling $\Omega_{\rm GW} \propto (H_{\rm bounce}/\Mpl)^2$ with a detector gap of $10^{17}$.

\textbf{Barrier 14 (Perturbation Transparency):} See Section~\ref{sec:transparency}.

\section{The Perturbation-Transparency Theorem}
\label{sec:transparency}

\subsection{Statement}

In minimal Einstein-Cartan-Holst gravity with canonical scalar field matter, the Holst term is dynamically inert for both scalar and tensor perturbations at all orders. The Barbero-Immirzi parameter $\gamma$ is invisible in all perturbation observables.

\subsection{Proof (Scalar Sector)}

The argument proceeds in five steps:

\begin{enumerate}
\item \textbf{Zero spin density.} A canonical scalar field $\phi$ has spin density $S^\lambda{}_{\mu\nu} = 0$ (by definition of ``canonical''---no spinor indices, no intrinsic angular momentum).

\item \textbf{Zero torsion.} In Einstein-Cartan theory, the torsion tensor is algebraically determined by the spin density: $T^\lambda{}_{\mu\nu} = 8\pi G\, (S^\lambda{}_{\mu\nu} + \delta^\lambda_{[\mu} S_{\nu]} + \ldots)$. With $S = 0$, we have $T^\lambda{}_{\mu\nu} = 0$ at all perturbation orders.

\item \textbf{Connection reduces to Levi-Civita.} With $T = 0$, the full connection $\Gamma^\lambda{}_{\mu\nu} = \mathring{\Gamma}^\lambda{}_{\mu\nu}$ (Christoffel symbols only).

\item \textbf{Holst term becomes topological.} The Holst term $\frac{1}{2\gamma}\int \epsilon^{IJ}{}_{KL}\, e^K \wedge e^L \wedge F_{IJ}$ evaluated with the Levi-Civita connection reduces to the Pontryagin density $\int \sqrt{-g}\, \epsilon^{\mu\nu\rho\sigma} R_{\mu\nu\rho\sigma}\, d^4x$, which is a total derivative in 4D.

\item \textbf{No equations of motion.} A total derivative contributes zero to the variational equations at all orders in perturbation theory.
\end{enumerate}

\subsection{Extension to Tensor Sector}

The same five-step argument applies to tensor perturbations: with $T = 0$, the tensor mode equation remains the standard GR form $h''_{ij} + 2\mathcal{H}h'_{ij} + k^2 h_{ij} = 0$, with no parity splitting ($\Delta v \equiv v_R - v_L = 0$ exactly). This was independently confirmed by the earlier Branch H parity-tensor analysis (Barrier 8).

\subsection{Implications}

ECH gravity is a \emph{perturbation-transparent UV regulator}: it resolves the singularity at Planck density through the $\rho^2$ modification of the Friedmann equation, but leaves zero fingerprint on the primordial perturbations. The Barbero-Immirzi parameter $\gamma$, which sets the bounce density $\rho_{\rm crit}$ and plays a central role in loop quantum gravity, is completely invisible in all scalar and tensor perturbation observables computed with canonical scalar field matter.

\section{The Hybrid Dark-Energy Loophole}

We considered the phenomenological route of appending late-time dynamical-dark-energy freedom (e.g., CPL $w_0w_a$ parametrization) to the bounce model. This was explored across 7 disguised forms:

\begin{enumerate}
\item Direct $w_0w_a$ addition to the background
\item Quintessence scalar with bounce-motivated initial conditions
\item Curvaton-derived late-time potential
\item Vacuum energy from cyclic boundary conditions
\item Torsion-induced effective $w(z)$
\item Holst-term residual as effective DE
\item ALP rolling as late-time acceleration
\end{enumerate}

All 7 forms were rejected on the following grounds: adding $w_0w_a$ to a bounce model produces the same fit improvement as adding $w_0w_a$ to $\Lambda$CDM, with no additional theoretical content from the bounce. None of the 236{,}622 MCMC posterior samples in this program used $w_0$ or $w_a$ as free parameters. The loophole was explored theoretically but never implemented computationally, because its implementation would not constitute first-principles evidence for bounce cosmology.

\section{Consequences for the Observable Program}

The 14 barriers, together with the perturbation-transparency theorem, establish that:

\begin{itemize}
\item ECH provides the \emph{background} for the bounce (singularity resolution at $\rho_{\rm crit}$)
\item ECH does \emph{not} provide distinctive perturbation-level observables
\item The observable science case for bouncing cosmology rests on \emph{generic}, mechanism-independent predictions
\item The strongest such prediction is $\fnl^{\rm local} = -35/8$~\citep{Cai:2009fn}, testable by SPHEREx at $4$--$6\sigma$ significance~\citep{Heinrich:2023}
\end{itemize}

This is not a failure of the framework---it is a sharp scientific conclusion. The perturbation-transparency of ECH means that the observable predictions of a matter bounce are \emph{more robust}, not less: they do not depend on the choice of bounce mechanism (ECH, LQC, or other UV completion).

\section{Conclusion}

We have presented a systematic catalog of 14 structural barriers that close all standard routes from nonsingular bounces to late-time dark energy or distinctive perturbation-level observables within the ECH framework. The central result is the perturbation-transparency theorem: minimal ECH is dynamically inert for scalar and tensor perturbations when the matter content is a canonical scalar field. The Holst term reduces to a topological invariant, and the Barbero-Immirzi parameter vanishes from all perturbation equations. The hybrid dark-energy loophole (appending phenomenological $w_0w_a$ freedom) was explored across 7 forms and rejected as non-derivational. The observable program for bouncing cosmology therefore targets generic, mechanism-independent predictions---principally $\fnl = -35/8$---rather than framework-specific signatures.

\section*{Acknowledgments}

The author thanks the developers of NumPy, SciPy, Matplotlib, mpmath, SymPy, and Tectonic. Computations were performed on consumer hardware. The author acknowledges helpful interactions with AI research assistants during the systematic barrier-cataloging and perturbation-gate phases.

\bibliography{focused_paper_refs}

\end{document}
